\cdcpartie{Conditions de réalisation}

\cdcsection{Enveloppe budgétaire}

Cette section précise les ressources que l’équipe compte mobiliser. Les montants dépendent des fournisseurs retenus et ne sont donc pas chiffrés à ce stade.

\cdcsubsection{Ressources humaines}

Le projet est réalisé par une équipe de cinq personnes. Chaque membre assure un rôle principal et participe aux revues collectives des livrables.

\begin{cdcitems}
  \item \textbf{Chef de projet.} Planification, coordination de l’équipe, suivi des jalons, relation avec l’encadrant et rédaction des procès-verbaux de réunion.
  \item \textbf{Analyste fonctionnel.} Recueil et formalisation des besoins, cahier des charges, modélisation des cas d’utilisation et des règles de gestion, maquettes.
  \item \textbf{Développeur front-end.} Interfaces web adaptées au smartphone, bilinguisme français et arabe, intégration de la carte.
  \item \textbf{Développeur back-end.} Base de données, logique métier des réparations, des paiements et des commissions, intégration de WhatsApp et du prestataire de paiement.
  \item \textbf{Responsable qualité et tests.} Plan de tests, tests automatisés, cahier de recette et suivi des anomalies.
\end{cdcitems}

\cdcsubsection{Ressources matérielles et logicielles}

L’équipe dispose déjà de ses postes de travail et utilise des outils gratuits pour le développement, la gestion de version et le suivi des tâches. Les autres ressources se répartissent en deux catégories.

\cdcintertitre{Coûts récurrents d’exploitation}

\begin{cdcitems}
  \item Hébergement web et base de données.
  \item Nom de domaine et certificat HTTPS.
  \item Stockage des photos de l’état initial, selon le volume.
\end{cdcitems}

\cdcintertitre{Coûts liés à l’usage}

\begin{cdcitems}
  \item API WhatsApp Business, facturée selon les messages envoyés.
  \item Service de cartographie, facturé selon le nombre de requêtes.
  \item Prestataire de paiement, qui prélève des frais sur chaque transaction.
\end{cdcitems}

\cdcintertitre{Principes budgétaires}

Pendant le développement, l’équipe utilise des offres gratuites et des environnements de test, en particulier le mode test du prestataire de paiement. Les coûts récurrents ne s’appliquent donc qu’à la mise en production. Les coûts liés à l’usage sont suivis par boutique, afin de pouvoir être pris en compte dans le modèle de commissions.
